\documentclass[12pt,a4paper,openany]{report}
\usepackage{tgcommon}

\title{Manual Didático para Defesa}
\author{João Matheus Del Vecchio França Barbosa}
\date{2026}

\begin{document}
\tgtitlepage
  {João Matheus Del Vecchio França Barbosa}
  {Do Dado Bruto ao Transformer}
  {Manual didático passo a passo para compreender e defender o trabalho}
  {2026}

\pagenumbering{roman}
\chapter*{Como usar este manual}

Este texto assume que o leitor começa sem domínio de microestrutura, processos
pontuais ou redes neurais. Cada capítulo responde quatro perguntas: o que é, por que
é necessário, como foi implementado e o que uma banca pode perguntar. Não memorize
equações isoladas; aprenda a narrativa causal:

\begin{quote}
mercado atualiza bid/ask \(\rightarrow\) mid-price muda \(\rightarrow\) surge evento
\(\rightarrow\) eventos se agrupam \(\rightarrow\) Hawkes modela agrupamento
\(\rightarrow\) MLE estima parâmetros \(\rightarrow\) rede aprende a imitar a
estimação rapidamente \(\rightarrow\) Transformer testa se atenção ajuda.
\end{quote}

\tableofcontents
\clearpage
\pagenumbering{arabic}

\chapter{A história inteira em linguagem simples}

Imagine que você observa uma fila. Pessoas não chegam uniformemente: às vezes ninguém
chega; às vezes chega um grupo. Um processo de Poisson diria que cada chegada não muda
a chance da próxima. Um processo de Hawkes diz que uma chegada pode aumentar
temporariamente a chance de outra.

No mercado, o ``cliente chegando'' é uma alteração de preço. Uma alteração pode
estimular outras porque algoritmos reagem, ordens são executadas, participantes
atualizam expectativas e liquidez muda. Hawkes resume essa dinâmica por quatro números:
\begin{itemize}
  \item \(\mu\): atividade que existiria sem eventos recentes;
  \item \(\alpha_{\mathrm{self}}\): quanto uma alta estimula altas, ou baixa estimula baixas;
  \item \(\alpha_{\mathrm{cross}}\): quanto alta estimula baixa, e vice-versa;
  \item \(\beta\): quão rápido o efeito desaparece.
\end{itemize}

MLE encontra os números tentando várias combinações e escolhendo a que melhor explica
os eventos. É relativamente lento. A proposta de Lee é simular muitos exemplos para
os quais os números verdadeiros são conhecidos e treinar uma LSTM: dada uma sequência,
ela cospe os quatro números em uma única passagem. Nosso trabalho reproduz isso em
cripto e testa Transformer.

\chapter{Passo 1: compreender o mercado observado}

\section{Bid, ask e spread}

\textbf{Bid} é o maior preço que alguém oferece para comprar. \textbf{Ask} é o menor
preço que alguém aceita para vender. Normalmente,
\[
 p_{\mathrm{bid}}\le p_{\mathrm{ask}}.
\]
A diferença é o spread. Exemplo: bid 41.734,90 e ask 41.735,00. Spread: 0,10.

\section{Mid-price}

O mid-price é o centro:
\[
 m=\frac{41.734{,}90+41.735{,}00}{2}=41.734{,}95.
\]
Ele não é necessariamente preço negociado. É uma medida do centro da melhor cotação,
menos dependente de saber se o último negócio bateu no bid ou no ask.

\section{Por que bookTicker}

\dataset{bookTicker} registra alterações do melhor bid e ask. \dataset{aggTrades}
registra negócios agregados. Para reproduzir o método, precisamos do primeiro. Trades
servem como validação de atividade, mas não recuperam com exatidão o melhor bid/ask.

\section{Por que futuros e não spot}

O acervo público histórico da Binance oferece \dataset{bookTicker} para futuros USD-M
no período necessário. O caminho oficial de spot retornou ausência. Isso não significa
que spot não tenha cotações ao vivo; significa que a fonte histórica gratuita escolhida
não oferece o artefato necessário. A escolha é metodológica, não preferência comercial.

\section{Perguntas da banca}

\begin{description}
  \item[Por que não usar último negócio?] Porque ele alterna entre lados do livro e
  adiciona bid--ask bounce. Mid-price representa melhor o centro da cotação.
  \item[Por que Binance?] Disponibilidade pública, granularidade e liquidez; não porque
  a Binance represente todo o mercado.
  \item[Por que perpétuos?] Porque o acervo histórico oferece top-of-book para USD-M.
  \item[Isso cria viés?] Sim. Resultados são condicionados à exchange e ao instrumento.
\end{description}

\chapter{Passo 2: transformar milhões de linhas em eventos}

\section{Por que 100 ms}

Dados chegam em milissegundos e contêm enorme atividade de atualização. Lee observa o
mid-price em intervalos de 0,1 s para estudar alta frequência, não ultra-alta frequência.
Em cada bin, usamos o último estado conhecido.

\section{Algoritmo mental}

\begin{enumerate}
  \item pegue atualizações entre 12:00:00,000 e 12:00:00,099;
  \item mantenha a última;
  \item calcule o mid-price;
  \item compare com o bin anterior;
  \item se aumentou, evento 1; se diminuiu, evento 2; se igual, nada;
  \item registre quanto tempo passou desde o evento anterior;
  \item divida o tamanho da mudança pelo tick size.
\end{enumerate}

\section{Interchegada}

Se eventos ocorreram em 3,8 s, 3,9 s e 4,1 s, interchegadas são 3,8 s a partir do
início, depois 0,1 s e 0,2 s. Redes recebem interchegada porque ela descreve ritmo.

\section{Tick}

Tick é o menor incremento permitido. Se BTC tem tick 0,1 e mid-price muda 2,5,
o salto tem 25 ticks. Usar ticks torna o tamanho comparável à estrutura do instrumento.

\section{Por que chunks}

Um dia pode ter 1,7 GB descompactado. Carregar tudo exigiria múltiplos gigabytes de RAM.
O programa lê um milhão de linhas por vez. Guarda o último bin incompleto para combiná-lo
com o próximo chunk; sem isso, uma fronteira artificial poderia criar estado errado.

\section{Contagens obtidas}

\begin{table}[ht]
\centering
\caption{Resultado da transformação principal.}
\begin{tabular}{lrr}
\toprule
Símbolo & Eventos & Janelas de 1.000 eventos\\
\midrule
BTCUSDT & 887.803 & 880\\
ETHUSDT & 939.101 & 933\\
\bottomrule
\end{tabular}
\end{table}

\section{Perguntas da banca}

\begin{description}
  \item[Por que manter a última atualização?] Ela representa o estado observável no
  fim do intervalo.
  \item[Por que eliminar preços iguais?] O modelo representa eventos de mudança.
  \item[Por que aparecem timestamps repetidos?] Várias atualizações podem compartilhar
  resolução temporal; reduzimos ao último estado antes da grade.
  \item[O filtro perde informação?] Sim, deliberadamente. Ele reduz ruído e alinha o
  estudo à escala do paper.
\end{description}

\chapter{Passo 3: processos pontuais}

\section{Do tempo regular ao tempo de eventos}

Uma série comum registra preço a cada segundo. Processo pontual registra os próprios
instantes em que algo aconteceu. A sequência é
\[
(t_1,c_1),(t_2,c_2),\ldots,
\]
em que \(c_k=1\) ou 2.

\section{Contagem}

\[
N_i(t)=\text{número de eventos do tipo }i\text{ até }t.
\]
Se ocorreram três altas até 10 s, \(N_1(10)=3\).

\section{Intensidade}

Intensidade é taxa condicional. Se \(\lambda_1(t)=5\) eventos/s, em intervalo pequeno
de 0,01 s esperamos aproximadamente 0,05 eventos de tipo 1. Não é dizer que ``ocorrerá
0,05 evento''; é uma média/probabilidade infinitesimal.

\chapter{Passo 4: Poisson e sua insuficiência}

No Poisson homogêneo:
\[
\lambda(t)=\lambda_0.
\]
O passado não importa. Interchegadas são exponenciais. Se dados exibem excesso de
intervalos curtos e rajadas, Poisson subestima agrupamento. Hawkes permite que cada
evento aumente a taxa futura.

\chapter{Passo 5: processo de Hawkes}

\section{Equação}

\[
\lambda_i(t)=\mu+\sum_j\sum_{t_k^j<t}\alpha_{ij}e^{-\beta(t-t_k^j)}.
\]
Leia da esquerda para direita:
\begin{enumerate}
  \item intensidade atual do tipo \(i\);
  \item começa em \(\mu\);
  \item soma influência de todos os eventos passados;
  \item influência inicial é \(\alpha_{ij}\);
  \item influência decai exponencialmente com distância temporal.
\end{enumerate}

\section{Self e cross}

Na matriz simétrica,
\[
\alpha=
\begin{bmatrix}\alpha_s&\alpha_c\\\alpha_c&\alpha_s\end{bmatrix}.
\]
\(\alpha_s\) mede persistência direcional do tipo de evento. \(\alpha_c\) mede reversão
ou reação entre tipos. A interpretação econômica deve ser cuidadosa: parâmetro não
prova causalidade comportamental.

\section{Beta e memória}

A meia-vida aproximada do kernel é
\[
t_{1/2}=\frac{\log 2}{\beta}.
\]
Beta grande: memória curta. Beta pequeno: efeito persiste.

\section{Razão de ramificação}

\[
g=\frac{\alpha_s+\alpha_c}{\beta}.
\]
Se \(g=0,8\), cada evento produz, em média, 0,8 filhos diretos. Filhos produzem netos;
o cluster esperado cresce, mas permanece finito. Se \(g\ge1\), a cascata pode não morrer.

\section{Perguntas da banca}

\begin{description}
  \item[Por que kernel exponencial?] É parcimonioso, possui recursão eficiente e segue
  o artigo-base.
  \item[Por que simétrico?] Reduz dimensão e reproduz Lee; não afirma que o mercado é
  perfeitamente simétrico.
  \item[Por que \(g<1\)?] Para estabilidade/estacionariedade.
  \item[Por que cross quase zero no MLE atual?] É indício, não conclusão; pode ser dado,
  identificação, otimização ou especificação e requer diagnóstico.
\end{description}

\chapter{Passo 6: simular Hawkes}

\section{Por que simular}

Em dados reais, os parâmetros verdadeiros são desconhecidos. Em simulação, escolhemos
\(\theta\), geramos sequência e perguntamos se o estimador recupera \(\theta\). Isso cria
treino supervisionado sem rótulo humano.

\section{Thinning de Ogata}

\begin{enumerate}
  \item calcule um limite superior \(M\) para intensidade total;
  \item sorteie espera \(w\sim\operatorname{Exp}(M)\);
  \item avance tempo e deixe excitações decaírem;
  \item calcule intensidade real \(\lambda_{\mathrm{tot}}\);
  \item aceite com probabilidade \(\lambda_{\mathrm{tot}}/M\);
  \item se aceito, sorteie tipo proporcionalmente a \(\lambda_1,\lambda_2\);
  \item adicione excitação e repita.
\end{enumerate}

O algoritmo rejeita alguns candidatos, exatamente como selecionar pontos de um processo
mais intenso para produzir o processo desejado \citep{ogata1981}.

\chapter{Passo 7: máxima verossimilhança}

\section{Intuição}

MLE procura parâmetros que tornem o observado plausível. A função é
\[
\ell=\underbrace{\sum_k\log\lambda_{c_k}(t_k)}_{\text{recompensa eventos}}
-\underbrace{\sum_i\int_0^T\lambda_i(u)\,du}_{\text{penaliza intensidade excessiva}}.
\]
Sem compensador, bastaria elevar todas as intensidades.

\section{Otimização}

L-BFGS-B procura máximo numericamente. Como pode depender da inicialização, usamos
múltiplos pontos. O código converte problema em minimização de \(-\ell\).

\section{Restrições}

\[
\mu=\softplus(z_1),\quad\alpha_s=\softplus(z_2),\quad
\alpha_c=\softplus(z_3),\quad d=\softplus(z_4),
\]
\[
\beta=\alpha_s+\alpha_c+d.
\]
Softplus, \(\log(1+e^z)\), é sempre positiva e diferenciável.

\section{Diagnóstico histórico, não final}

Em uma verificação inicial, doze ajustes representativos convergiram. Esse registro
histórico apenas confirma que o pipeline executou com janelas de 1.000 eventos; ele
não constitui medição final de erro ou desempenho. As métricas finais vêm
exclusivamente de \code{thesis/latex/final_results.tex}.

\chapter{Passo 8: LSTM de Lee}

\section{Entrada}

Cada passo tem dois números: \(\log(1+\Delta t)\) e tipo codificado como \(+1/-1\).
O log comprime interchegadas muito grandes.

\section{Memória recorrente}

A LSTM lê um passo, atualiza memória, lê o seguinte e assim por diante. Duas camadas
com 12 unidades produzem estado final. Uma camada linear produz quatro logits, depois
transformados em parâmetros válidos.

\section{Inferência amortizada}

MLE otimiza novamente para cada janela. LSTM treina uma vez, depois apenas executa
multiplicações. Treino pode ser caro; inferência tende a ser barata. ``Mais rápida''
deve ser comparado no mesmo hardware, lote e protocolo.

\section{O que é MSE}

\[
\MSE=\frac1n\sum_{r=1}^{n}\|\hat\theta_r-\theta_r\|^2.
\]
Erros grandes pesam quadraticamente. Também precisamos métricas por parâmetro porque
escalas diferentes podem dominar a média.

\chapter{Passo 9: Transformer}

\section{Atenção}

Cada evento cria query, key e value. Produto query--key indica relevância. Softmax
transforma escores em pesos. Value combina informação. Múltiplas cabeças podem aprender
padrões diferentes.

\section{Posição}

Sem encoding, Transformer não sabe ordem. Adicionamos senoides dependentes da posição.
Uma extensão futura pode usar encoding do tempo acumulado, pois eventos são irregulares.

\section{O caminho dos tensores}

Começamos com uma caixa de formato \([B,L,2]\): \(B\) sequências, \(L\) eventos e
dois atributos por evento. Uma camada linear transforma cada par em 32 números, dando
\([B,L,32]\). Somamos posição. Quatro cabeças recebem oito números por posição cada
uma e aprendem relações diferentes. O bloco feed-forward expande cada posição de 32
para 128 números e volta para 32; isso ocorre duas vezes, pois há duas camadas encoder.

Não são quatro redes independentes: as cabeças são partes paralelas da mesma camada.
Query pergunta ``o que procuro?'', key descreve ``o que ofereço?'' e value contém
``qual informação entrego?''. O produto query--key, seguido de softmax, decide quanto
de cada value será misturado na posição atual.

\section{Pooling}

Encoder produz vetor por posição. A média resume sequência em um vetor, então a mesma
cabeça estável gera parâmetros.

A média muda \([B,L,32]\) para \([B,32]\). A cabeça produz quatro números positivos:
\(\mu,\alpha_s,\alpha_c,d\). O quarto parâmetro apresentado é construído como
\(\beta=\alpha_s+\alpha_c+d\). Assim
\((\alpha_s+\alpha_c)/\beta<1\) automaticamente; não dependemos de consertar uma
previsão instável depois.

\section{Comparação justa}

Transformer e LSTM devem compartilhar dataset, splits, targets, loss e orçamento
de busca. Caso contrário, não sabemos se ganho veio da arquitetura ou de tratamento
preferencial.

\chapter{Passo 10: interpretar os resultados sem exagerar}

\section{Diagnósticos históricos e resultados finais}

Os ensaios reduzidos usados durante o desenvolvimento eram apenas verificações de
software. Seus números estão obsoletos e não devem ser usados como evidência
científica. As métricas finais sintéticas e empíricas, inclusive erro e latência,
são reportadas somente pelo fragmento
\code{thesis/latex/final_results.tex}, incorporado neste manual no Passo 11.

\section{O novo protocolo executável}

O protocolo atual usa exclusivamente sequências de 1.000 eventos. A configuração
local usa 48/16/16 trajetórias e duas épocas apenas para encontrar erros de software.
A configuração final usa 75.000 trajetórias para treino, 10.000 para validação,
15.000 para teste, até 300 épocas e cinco seeds. O comando \code{make smoke} valida
o pipeline, mas não produz métricas científicas finais; os comandos finais são
separados porque podem consumir horas ou dias.

Nos dados reais, todas as janelas também contêm 1.000 eventos. O treinamento só pode
enxergar 15--22 de janeiro. A escolha usa 23--25. O teste 26--28 exige abertura
explícita. Outubro usa o checkpoint congelado, sem treinar novamente. Se a banca
perguntar por que, a resposta é: permitir que o modelo aprenda com o teste transforma
medição de generalização em memorização.

\section{Frases corretas e incorretas}

\begin{description}
  \item[Correta:] ``O protocolo final compara os modelos em sequências de 1.000 eventos.''
  \item[Incorreta:] ``Há também uma execução final com outro comprimento de sequência.''
  \item[Correta:] ``As métricas finais vêm de
  \code{thesis/latex/final_results.tex}.''
  \item[Incorreta:] ``Os números das verificações de software são resultados finais.''
  \item[Correta:] ``A fonte escolhida cobre as janelas estudadas.''
  \item[Incorreta:] ``Os dados representam todo o mercado de cripto.''
\end{description}

\chapter{Passo 11: reprodutibilidade}

\section{Manifesto}

Cada ZIP recebe URL, tamanho e SHA-256. SHA-256 funciona como impressão digital:
se um byte mudar, hash muda com altíssima probabilidade.

\section{Parquet}

CSV é volumoso e sem tipos fortes. Parquet é colunar, comprimido e eficiente para ler
somente colunas necessárias. O CSV temporário é apagado após conversão.

\section{Seeds}

Seed torna sorteio reproduzível. Não elimina aleatoriedade científica; permite repetir
uma realização. Resultado final precisa de várias seeds.

\section{Git}

Código e configurações entram no Git; dados pesados, checkpoints e artefatos não.
O manifesto permite reconstruí-los. Um commit identifica exatamente a versão do código.

\section{Fonte única das métricas finais}

As métricas finais não são digitadas manualmente neste manual. Elas são geradas em
\code{thesis/latex/final_results.tex} a partir do resumo consolidado e incorporadas
abaixo.

\input{final_results.tex}

\chapter{Passo 12: como apresentar à banca}

\section{Roteiro de quinze minutos}

\begin{enumerate}
  \item 1 min: problema---clustering de mudanças de preço;
  \item 2 min: Hawkes e interpretação dos quatro parâmetros;
  \item 2 min: gargalo do MLE e ideia de Lee;
  \item 2 min: gap---cripto e atenção;
  \item 2 min: dados e transformação;
  \item 2 min: desenho experimental sem vazamento;
  \item 2 min: métricas finais do fragmento gerado;
  \item 1 min: limitações;
  \item 1 min: próximos passos e contribuição.
\end{enumerate}

\section{Perguntas difíceis}

\begin{description}
  \item[Por que usar MLE como rótulo real?] Porque parâmetro verdadeiro não existe em
  dados reais; imitar MLE é objetivo operacional. Mas herda seus vieses.
  \item[Por que não Neural Hawkes?] Porque o objetivo é estimar parâmetros
  interpretáveis do Hawkes clássico, não substituir sua intensidade por caixa-preta.
  \item[Como evitar vazamento?] Split por blocos temporais antes de gerar janelas.
  \item[Por que duas semanas?] Volume já gera milhões de eventos e mantém execução
  administrável; robustez é testada em semana separada.
  \item[Por que usar apenas 1.000 eventos?] Porque esse é o comprimento único fixado
  pelo protocolo vigente. Resultados com outros comprimentos não integram a execução
  final nem as comparações reportadas.
  \item[Qual contribuição se Transformer perder?] Reprodução, dataset/pipeline pública
  e evidência controlada de que atenção não traz ganho nesse regime.
\end{description}

\chapter{Passo 13: extensões após o protocolo final}

\begin{enumerate}
  \item ampliar o MLE empírico nas janelas de 1.000 eventos;
  \item abrir o teste empírico fechado uma única vez;
  \item avaliar outubro sem tuning nem retreinamento;
  \item diagnosticar resíduos do Hawkes;
  \item calcular e comparar volatilidades;
  \item revisar escrita, referências e normas ITA.
\end{enumerate}

\appendix
\chapter{Glossário}

\begin{description}
  \item[Ask:] menor preço ofertado por vendedor.
  \item[Bid:] maior preço ofertado por comprador.
  \item[Evento:] mudança de mid-price após filtro.
  \item[Inferência amortizada:] custo de treino reutilizado em inferências futuras.
  \item[Intensidade:] taxa condicional instantânea de eventos.
  \item[Kernel:] forma temporal do efeito de um evento.
  \item[MLE:] estimativa por máxima verossimilhança.
  \item[Mid-price:] média entre bid e ask.
  \item[Self-attention:] mecanismo que relaciona posições da mesma sequência.
  \item[Split:] partição de treino, validação ou teste.
  \item[Tick:] menor incremento de preço permitido.
\end{description}

\chapter{Comandos}

\begin{lstlisting}[language=bash]
uv sync --python 3.13 --extra dev
uv run python scripts/build_primary_dataset.py --workers 4
uv run python scripts/index_event_windows.py
make smoke
make synthetic-final
make train-final
make calibration
make empirical-test
uv run python scripts/generate_final_report_fragment.py
uv run ruff check .
uv run mypy src/hawkes_crypto
uv run pytest
\end{lstlisting}

\bibliographystyle{plainnat}
\bibliography{references}
\end{document}

